% Nháp Mục 3 (Method) theo trọng tâm mới (Review-2026-11-17, TQĐ28; Bảng 7 của review). Mẫu LNCS, không dùng "we".
% Không chứa kết quả của pha xác nhận. Số của cơ chế lấy từ docs/taimma/disruption.csv (mô phỏng của GVHD, 200 lần lặp mỗi ô,
% p0 = 5000, k = 10) và T1/history.csv của S1 (tag taimma-2027-v1-results). Bản cũ draft_03_method.tex (v1) giữ để đối chiếu.
% Cần trong preamble: \usepackage{amsmath,amssymb} \usepackage{algorithm} \usepackage[noend]{algpseudocode}

\section{Method}\label{sec:method}

\subsection{Problem and fitness}\label{sec:problem}

Let $X\in\mathbb{R}^{n\times p_0}$ and $y\in\{0,1\}^n$ be the training part of one outer fold after a label-free variance
prefilter that keeps $p_0=5{,}000$ genes. Given a budget $k$, a wrapper searches a panel $S\subset\{1,\dots,p_0\}$ with
$|S|=k$ that maximises
\begin{equation}\label{eq:fitness}
  f(S) = -\frac{1}{3}\sum_{q=1}^{3}\mathrm{LL}_q(S),
\end{equation}
where $\mathrm{LL}_q$ is the log-loss on the validation part of inner fold $q$ of an $L_2$-regularised logistic regression
($C=1$, balanced class weights) fitted on the standardised genes of $S$; probabilities are clipped to $[10^{-6},1-10^{-6}]$.
The three inner folds are fixed and shared by all methods on the same outer fold. The outer test part is never read during
the search. Every method receives the same budget of $P=30$ individuals and $T=50$ generations, so at most
$P(T+1)+1=1{,}531$ distinct panels are evaluated.

\subsection{Random keys and band re-encoding (v1)}\label{sec:v1}

An individual is a key vector $x\in(0,1)^{p_0}$ and its panel is $S(x)=\mathrm{top}_k(x)$, the indices of the $k$ largest
keys~\cite{bean1994}. Every individual is feasible. The first version of the method, MPQDA-DE v1, re-encodes every vector
that enters the population into two bands while keeping its panel and the order of keys within each group:
\begin{equation}\label{eq:band}
  \mathrm{band}_\varepsilon(x)_j=
  \begin{cases}
    \tfrac12+\tfrac12\,\dfrac{r_j+1/2}{k}, & j\in S(x),\\[4pt]
    \varepsilon\,\dfrac{r_j+1/2}{p_0-k}, & j\notin S(x),
  \end{cases}
\end{equation}
with $\varepsilon=1/2$ and $r_j$ the stable rank of $x_j$ within its group. v1 draws the leader of each mutation from an archive
of near-optimal panels by a roulette on novelty and fitness, controls the scale factor $F$ of each individual by Q-learning,
initialises the keys with the logistic map and returns a consensus panel computed from the archive. Its evaluation did not
support the expected advantages: the search did not reach a better inner fitness than DE/rand/1/bin, and neither the
Q-learning control nor the chaotic initialisation had a measurable effect (Section~\ref{sec:results-v1}).

\subsection{Panel disruption}\label{sec:disruption}

Genes outside the panel are under no selection pressure, so their keys keep a random order inside the lower band
$(0,\varepsilon)$. For a gene $j$ outside the panels of the base vector and of both difference vectors, the mutant key is
\begin{equation}\label{eq:mutant}
  v_j = x_{Lj} + F\,(x_{r_1 j}-x_{r_2 j}) \in \big(-F\varepsilon,\ \varepsilon(1+F)\big).
\end{equation}
With $\varepsilon=1/2$ this interval overlaps the panel band $(1/2,1)$, so the difference vector lifts out-of-panel genes above
the low-ranked genes of the panel. A trial panel then carries genes that belong to no parent. A simulation of the mutation and
binomial crossover operators alone ($p_0=5{,}000$, $k=10$, $CR=0.9$, 200 repetitions per setting) quantifies the effect: at
$F=0.5$, 34--36\,\% of the genes of a trial panel are random genes and the trial keeps 47--49\,\% of the genes of its leader;
the share of random genes grows from 18--22\,\% at $F=0.3$ to 68--69\,\% at $F=0.9$.
% source: docs/taimma/disruption.csv, CR = 0.9, eps = 0.5, share ∈ {0, 0.5}
On the three GPL570 cohorts, the trial panels of v1 kept on average 44\,\% of the genes of their leader in the first
generation and 53\,\% in the last one.
% source: T1/<cohort>/history.csv (S1), method = mpqda_de, leader_overlap, t = 1 and t = 50
The scale factor therefore sets the amount of random perturbation rather than the direction of the move, which is consistent
with the absence of an effect of the Q-learning control of $F$.

\subsection{MPQDA-DE v2}\label{sec:v2}

MPQDA-DE v2 keeps the representation, the fitness, the archive, the consensus output and the budget of v1, and changes three
parts. The constants below were fixed from the argument of Section~\ref{sec:disruption}, not tuned on data.

\paragraph{Gapped re-encoding.} The lower band of~\eqref{eq:band} is compressed to $\varepsilon=0.01$. The bound
in~\eqref{eq:mutant} becomes $\varepsilon(1+F)\le 0.02$, far below the panel band, so mutation and crossover only exchange genes
between the panels of the leader, of the added vector and of the target. In the simulation above the share of random genes in
a trial panel falls to 0 and the trial keeps 84--87\,\% of the genes of its leader at $F=0.5$.
% source: docs/taimma/disruption.csv, CR = 0.9, eps = 0.01, F = 0.5

\paragraph{New genes.} New genes enter only through an explicit step: after crossover, each gene of $S(u_i)$ is replaced
independently with probability $1/k$ by a gene drawn uniformly outside the panel, by swapping the two keys. A trial thus
receives one new gene on average, as in the mutation of a genetic algorithm.

\paragraph{Archive tournament.} The leader $L_i$ is the best of three archive members drawn uniformly with replacement, ranked
by fitness; the vector $x_{r_1}$ of the difference is replaced by another archive member chosen in the same way. $F=0.5$ and
$CR=0.9$ are fixed, and the keys are initialised uniformly; the Q-learning control and the chaotic initialisation of v1 are
removed.

\begin{algorithm}[t]
\caption{MPQDA-DE v2 for one outer fold and one budget $k$}\label{alg:v2}
\begin{algorithmic}[1]
\Require training part $(X,y)$ after the prefilter, budget $k$, seed; $P=30$, $T=50$, $\varepsilon=0.01$, $F=0.5$, $CR=0.9$
\State $x_i\sim U(0,1)^{p_0}$; $S_i\gets\mathrm{top}_k(x_i)$; $x_i\gets\mathrm{band}_\varepsilon(x_i)$; $f_i\gets f(S_i)$; offer $(S_i,f_i,x_i)$ to the archive $A$
\For{$t=1,\dots,T$}
  \State for each $i$: $L_i\gets$ best of 3 members of $A$; $a_i\gets$ best of 3 members of $A$; $r_2\neq i$ uniform in the population
  \State for each $i$: $v_i\gets\mathrm{repair}\big(x_{L_i}+F(x_{a_i}-x_{r_2}),\,x_i\big)$; $u_i\gets$ binomial crossover of $x_i$ and $v_i$
  \State for each $i$: replace each gene of $\mathrm{top}_k(u_i)$ with probability $1/k$ by a gene outside the panel (swap keys)
  \State evaluate $f(\mathrm{top}_k(u_i))$ for all $i$ in one batch
  \State for each $i$ with $f(u_i)\ge f(x_i)$: $x_i\gets\mathrm{band}_\varepsilon(u_i)$; $f_i\gets f(u_i)$; offer $(\mathrm{top}_k(u_i),f_i,x_i)$ to $A$
\EndFor
\State $S_c\gets$ consensus of $A$; evaluate $f(S_c)$
\State \Return $S_c$ (primary output) and the best individual
\end{algorithmic}
\end{algorithm}

\paragraph{Consensus output.} The archive keeps at most 15 panels that differ by a Jaccard distance of at least $0.2$, as in v1.
At termination, members are weighted by fitness rank, $w_a=|A|-\mathrm{rank}(a)+1$, each gene receives
$\pi_j=\sum_a w_a\mathbb{1}[j\in S_a]/\sum_a w_a$, and the consensus panel $S_c$ holds the $k$ genes with the largest $\pi_j$.
Aggregating several near-optimal panels addresses the non-uniqueness of gene signatures~\cite{einDor2005}; the best
individual of the same run is reported for comparison.

\paragraph{Internal control.} DE-G is DE/rand/1/bin with the gapped re-encoding and the new-gene step but without the archive
and the tournament. Comparing DE-G with DE/rand/1/bin isolates the effect of the representation, and comparing v2 with DE-G
isolates the effect of the archive tournament.

\paragraph{Selection of the configuration.} Three candidates were compared on separate development data (three
patient-disjoint subsets of the TCGA LUAD/LUSC cohort and GSE42127) by a rule written before the runs: the configuration
above (fixed $F$ and $CR$, tournament of size 3), a variant with SHADE-type adaptation of $F$ and $CR$, and a variant with a
tournament of size 2. The first one had the smallest mean rank of inner fitness and was retained
(Section~\ref{sec:protocol}).

% Khoá BibTeX: bean1994, einDor2005 (doi:10.1093/bioinformatics/bth469); Reunanen 2003 và Cawley, Talbot 2010 dùng ở Mục 4–6.
% Viết ở thể bị động hoặc lấy phương pháp làm chủ ngữ; không có số của pha xác nhận trong mục này.
